% ============================================================================
% KAPITEL 15: Testbare Vorhersagen
% ============================================================================

\chapter{Testbare Vorhersagen}
\label{ch:testbare_vorhersagen}

\section{Einleitung}
\label{sec:vorh_einleitung}

Die WDBT+ macht eine Reihe spezifischer, von der Allgemeinen Relativitätstheorie und dem Standardmodell der Kosmologie abweichender Vorhersagen. Diese Vorhersagen sind prinzipiell experimentell überprüfbar – entweder mit bestehenden oder mit zukünftigen Technologien.

Die folgende Liste fasst die wichtigsten testbaren Vorhersagen zusammen, geordnet nach physikalischen Bereichen.

\section{Kosmologie und Galaxien}
\label{sec:vorh_kosmologie}

\subsection{Fraktale Skalengesetze}
\label{sub:vorh_fraktale_skalen}

Die fraktale Raumdimension \( D \approx 2,704 \) (siehe Kapitel~\ref{ch:fraktale_dimension}) führt zu gebrochenen Exponenten in kosmologischen Skalengesetzen:

\begin{itemize}
    \item Galaxien-Dichteprofil: \( \rho(r) \propto r^{-(3-D)} \approx r^{-0,296} \) (flacher als das NFW-Profil)
    \item Rotationskurven: \( v(r) \propto r^{(D-2)/2} \approx r^{0,352} \) (bei sehr großen Radien steigt die Kurve langsam an, fällt nicht ab)
    \item CMB-Korrelationen: \( C(\theta) \propto \theta^{-(3-D)} \approx \theta^{-0,296} \) mit einer festen, aus der Dodekaeder-Symmetrie folgenden Achse (\enquote{Achse des Bösen})
    \item Materieentstehung: \( \frac{dM}{d\ln r} \propto r^{3-D} \approx r^{0,296} \)
\end{itemize}

Die \enquote{Achse des Bösen} im CMB ist dabei eine direkte Konsequenz der Dodekaeder-Geometrie selbst. Ein Dodekaeder besitzt eine 5-zählige Symmetrie mit ausgezeichneten Achsen (Verbindungslinien zwischen gegenüberliegenden Flächenmittelpunkten oder Ecken). Eine raumfüllende Packung von Dodekaedern überträgt diese Vorzugsrichtung auf die großskalige Struktur des Universums. Die Achse des Bösen ist somit \emph{nicht} chiral – sie ist eine Richtung, die aus der Geometrie des Dodekaeders folgt. Die Chiralität des Raumes hingegen entsteht erst durch die Verdrehung zwischen den Iterationen der Fibonacci-Rekursion (siehe Abschnitt~\ref{sub:randbedingung_packung}) und ist ein davon unabhängiges Phänomen.

\subsection{Maximale Masse kompakter Objekte}
\label{sub:vorh_max_masse}

Die Theorie sagt eine natürliche Obergrenze für kompakte Objekte voraus (siehe Kapitel~\ref{ch:bec_objekte}):

\[
M_{\text{BEC}} \approx 3 \times 10^6 M_{\odot}
\]

Beobachtete supermassereiche Objekte mit Massen über dieser Grenze sind keine einzelnen kompakten Objekte, sondern Systeme aus einem BEC-Kern und einer massiven Umgebung.

\subsection{Doppelkern-Systeme}
\label{sub:vorh_doppelkern}

Bei Galaxienverschmelzungen sagt die Theorie voraus, dass Doppelkern-Systeme häufiger sind als in der ART vorhergesagt, da BEC-Objekte bei Überschreiten der maximalen Masse zerreißen können, anstatt zu verschmelzen.

\subsection{Hubble-Gesetz und Hubble-Spannung}
\label{sub:vorh_hubble}

Die Rotverschiebung ist keine Doppler-Folge der Expansion, sondern eine kumulative gravitative Wirkung (siehe Kapitel~\ref{ch:kosmologie}). Die effektive Hubble-Konstante ist daher skalenabhängig:

\[
H_{\text{eff}}(d) = H_0 \left( \frac{d}{d_0} \right)^{D-2} = H_0 \left( \frac{d}{d_0} \right)^{0,704}
\]

Die Hubble-Spannung – die Diskrepanz zwischen CMB- und lokalen Messungen – ist keine Inkonsistenz, sondern eine direkte Konsequenz der fraktalen Geometrie:

\[
\frac{H_{\text{local}}}{H_{\text{CMB}}} = \left( \frac{d_{\text{local}}}{d_{\text{CMB}}} \right)^{0,704}
\]

Mit den gemessenen Werten \( H_{\text{CMB}} \approx 67,4 \, \text{km/s/Mpc} \) und \( H_{\text{local}} \approx 73,5 \, \text{km/s/Mpc} \) ergibt sich:

\[
\frac{d_{\text{local}}}{d_{\text{CMB}}} = \left( \frac{73,5}{67,4} \right)^{1/0,704} \approx 1,129
\]

Die lokalen Messungen operieren also auf einer um etwa 12,9 \% \textbf{größeren} effektiven Skala als die CMB-Messungen. Da \( H_{\text{eff}}(d) \) mit der Skala \( d \) wächst (Exponent \( D-2 = 0,704 > 0 \)), ist der aus lokalen Messungen abgeleitete Wert \( H_{\text{local}} \) \textit{größer} als der CMB-Wert \( H_{\text{CMB}} \), was qualitativ mit den Beobachtungen übereinstimmt.

\subsection{Übersicht der testbaren Vorhersagen}
\label{sub:vorh_uebersicht}

Die folgende Tabelle fasst die wichtigsten Vorhersagen zusammen:

\begin{tabular}{|l|l|l|}
\hline
\textbf{Vorhersage} & \textbf{Abweichung von ART} & \textbf{Experiment} \\
\hline
Lichtablenkung & \(\Delta\phi = \Delta\phi_{\text{ART}}(1 + \alpha\lambda^2)\) & EHT, VLBI, LISA \\
Gravitationswellen-Dispersion & \(v_{\text{phase}} = c \left(1 + \beta \left(\frac{f}{f_0}\right)^{D-3}\right)\) & LISA, Einstein-Teleskop \\
Neutrinomasse & \(m_\nu \approx 0,1\,\text{eV}/c^2\) (drei Generationen) & KATRIN, JUNO, DUNE \\
Hubble-Skalenabhängigkeit & \(H_{\text{eff}}(d) = H_0 (d/d_0)^{0,704}\) & Supernovae, Quasare, GW \\
Maximale kompakte Masse & \(M_{\text{BEC}} \approx 3 \times 10^6 M_\odot\) & Galaxienzentren, EHT \\
Galaxien-Dichteprofil & \(\rho(r) \propto r^{-0,296}\) & EUCLID, Rubin (LSST) \\
Rotationskurven & \(v(r) \propto r^{0,352}\) (asymptotisch steigend) & Extrem breite Galaxien \\
\hline
\end{tabular}

Die benötigte relative Genauigkeit liegt für die meisten Experimente im Bereich \( 10^{-6} \) bis \( 10^{-2} \). Jede dieser Vorhersagen ist prinzipiell falsifizierbar.

\section{Gravitationsphysik}
\label{sec:vorh_gravitation}

\subsection{Frequenzabhängige Lichtablenkung}
\label{sub:vorh_lichtablenkung}

Die ART sagt eine frequenzunabhängige Lichtablenkung voraus. Die WDBT+ sagt dagegen eine frequenzabhängige Ablenkung voraus:

\[
\Delta\phi(\nu) = \Delta\phi_0 \left(1 + \alpha \frac{\nu_0}{\nu}\right) \propto \lambda^2
\]

Dabei ist \( \alpha \) eine dimensionslose Konstante, die aus der fraktalen Raumdimension \( D \) folgt (siehe Anhang~\ref{app:herleitungen}, Gleichung D.12). Dieser Effekt ist mit hochpräzisen Interferometern (z. B. LISA) oder spektral aufgelösten Messungen am Sonnenrand testbar.

\subsection{Gravitationswellen-Dispersion}
\label{sub:vorh_gw_dispersion}

Die fraktale Raumstruktur führt zu einer frequenzabhängigen Ausbreitungsgeschwindigkeit von Gravitationswellen (siehe Anhang~\ref{app:spektrale_dimension}). Aus der spektralen Dimension \( d_s = 2D/3 \) folgt die Dispersionsrelation \( \omega \propto k^{D/3} \), und daraus ergibt sich für die Phasengeschwindigkeit:

\[
v_{\text{phase}}(\omega) = \frac{\omega}{k} \propto \omega^{1 - D/3}
\]

Mit \( D \approx 2,704 \) ist \( 1 - D/3 \approx -0,099 \), also eine schwache Dispersion bei hohen Frequenzen. In der üblichen Form mit einer Referenzfrequenz:

\[
v_{\text{phase}}(\omega) = c \left(1 + \beta \left(\frac{\omega}{\omega_0}\right)^{1 - D/3} + \dots\right)
\]

Dies ist mit zukünftigen Detektoren wie LISA (\cite{LISA2017}) oder dem Einstein-Teleskop (\cite{EinsteinTeleskop2010}) testbar.

\subsection{Shapiro-Laufzeitverzögerung}
\label{sub:vorh_shapiro}

Die ART sagt eine frequenzunabhängige Laufzeitverzögerung voraus. Die WDBT+ sagt eine zusätzliche Frequenzabhängigkeit voraus:

\[
\Delta t_{\text{WG}} \propto \lambda^{-2}
\]

Dies ist mit Pulsar-Timing-Experimenten testbar.

\subsection{Spiralbahnen statt Ellipsen}
\label{sub:vorh_spiralbahnen}

Die DSTT sagt voraus, dass alle Bahnen logarithmische Spiralen sind, keine geschlossenen Ellipsen (siehe Kapitel~\ref{ch:dstt_weber_kraefte}). Die radiale Drift ist extrem langsam:

\[
\frac{\dot{r}}{r} \approx 10^{-11} \text{ pro Jahr}
\]

Für den Mond ergibt sich eine Exzentrizitätsabnahme von \( \dot{e} \approx -2,7 \times 10^{-11} \) pro Jahr. Dieser Effekt ist mit heutigen Methoden (Lunar Laser Ranging) nicht nachweisbar, aber prinzipiell testbar.

\section{Quanten- und Teilchenphysik}
\label{sec:vorh_quanten}

\subsection{Neutrinomasse}
\label{sub:vorh_neutrinomasse}

Die Theorie sagt eine Neutrinomasse von (siehe Kapitel~\ref{ch:neutrino})

\[
m_{\nu,i} = \frac{\hbar}{l_0 c} \left( \frac{l_0}{L_{Q,0}} \cdot \left(\frac{2}{3-i}\right)^{\frac{3-D}{2}} \right)^{\frac{3-D}{2}}
\]

mit \( L_{Q,0} \approx 2{,}0 \times 10^{46} \, \text{m} \) und \( i = 1,2,3 \) für die drei Generationen.

Die numerischen Werte sind:

\[
\begin{array}{c|c}
\text{Generation} & m_{\nu,i} \\
\hline
1 & 0,101 \, \text{eV}/c^2 \\
2 & 9,3 \times 10^{-4} \, \text{eV}/c^2 \\
3 & \ll 10^{-4} \, \text{eV}/c^2
\end{array}
\]

Die leichteste Generation hat also:

\[
\boxed{m_{\nu,1} \approx 0,1 \, \text{eV}/c^2}
\]

Dies ist mit Experimenten wie KATRIN (\cite{KATRIN2022}) oder dem neutrinolosen Doppelbeta-Zerfall testbar.

\subsection{Neutrino-Oszillationen}
\label{sub:vorh_neutrino_osz}

Die fraktale Dispersionsrelation mit der fundamentalen Skala \( L_{Q,0} \approx 2{,}0 \times 10^{46} \, \text{m} \) führt zu einer charakteristischen Skalierung der Oszillationslänge:

\[
L_{\text{osc}} \propto E^{1/(4-D)} = E^{0,775}
\]

Dies weicht von den Vorhersagen des Standardmodells (\( L_{\text{osc}} \propto E \)) ab und ist mit Neutrino-Oszillationsexperimenten (JUNO, DUNE, Hyper-Kamiokande) testbar (siehe auch Anhang~\ref{app:neutrino_oszillation}).

Für eine gegebene Energie \( E \) ist die Oszillationslänge um den Faktor

\[
\frac{L_{\text{osc, WDBT}}}{L_{\text{osc, SM}}} \propto E^{-0,225}
\]

gegenüber dem Standardmodell verändert – ein Effekt, der bei hohen Energien sichtbar wird.

\subsection{Drei Neutrino-Generationen}
\label{sub:vorh_drei_generationen}

Die drei Generationen sind Moden des Q-Feldes im fraktalen Raum mit unterschiedlichen Korrelationslängen (siehe Kapitel~\ref{ch:neutrino}):

\[
L_{Q,i} = L_{Q,0} \cdot \left(\frac{3-i}{2}\right)^{\frac{2}{3-D}}
\]

Die Massenunterschiede zwischen den Generationen sind:

\[
\Delta m_{12}^2 = m_{\nu,1}^2 - m_{\nu,2}^2 \approx 0,01 \, \text{eV}^2/c^4
\]

\[
\Delta m_{23}^2 = m_{\nu,2}^2 - m_{\nu,3}^2 \approx 10^{-3} \, \text{eV}^2/c^4
\]

Dies ist konsistent mit den beobachteten Oszillationsparametern (\( \Delta m_{12}^2 \sim 7,5 \times 10^{-5} \, \text{eV}^2 \), \( \Delta m_{23}^2 \sim 2,5 \times 10^{-3} \, \text{eV}^2 \)) innerhalb der Größenordnung.

\subsection{Modifizierter Lamb-Shift}
\label{sub:vorh_lamb_shift}

Die Wechselwirkung des Elektrons mit dem Quantenpotential des Vakuums führt zu einem Zusatzterm im Lamb-Shift:

\[
\Delta E_{\text{Lamb}}^{\text{WDBT}} = \Delta E_{\text{QED}} + \frac{e^2 \hbar}{4 \pi \epsilon_0 m_e^2 c^3} \langle r \rangle
\]

Dabei ist \( \langle r \rangle \) der Erwartungswert des Ortes des Elektrons im gebundenen Zustand. Dieser Term ist klein, aber prinzipiell messbar.

\subsection{Verbindung zur Hubble-Spannung}
\label{sub:vorh_neutrino_hubble}

Eine der bemerkenswertesten Vorhersagen der WDBT+ ist die Verbindung zwischen der Neutrinomasse und der Hubble-Spannung. Beide Phänomene folgen aus derselben fraktalen Geometrie:

\begin{itemize}
    \item Die Neutrinomasse wird durch die Korrelationslänge \( L_{Q,0} \approx 2{,}0 \times 10^{46} \, \text{m} \) bestimmt.
    \item Dieselbe Skala \( L_{Q,0} \) bestimmt die effektive Hubble-Konstante.
    \item Die Hubble-Spannung ist daher kein Rätsel, sondern eine experimentelle Bestätigung der fraktalen Struktur des Raumes.
\end{itemize}

Die konsistente Erklärung beider Phänomene aus einer einzigen fundamentalen Skala \( L_{Q,0} \) ist ein starkes Indiz für die Richtigkeit der Theorie.

\section{Plasmaphysik}
\label{sec:vorh_plasma}

\subsection{Konstanz der Sonnenmasse}
\label{sub:vorh_sonnenmasse}

Die WDBT+ sagt voraus, dass der Sonnenwind nicht aus der Sonne selbst stammt, sondern aus neu entstehender Materie (siehe Kapitel~\ref{ch:kosmologie}). Die Sonne verliert daher keine signifikante Masse. Präzise Messungen der Sonnenmasse (z. B. durch Satellitenbahnen oder das MESSENGER-Experiment) über Zeiträume von Jahrzehnten sollten diesen Effekt bestätigen – oder die Theorie falsifizieren.

\subsection{Fraktale Skalierung im Sonnenwind}
\label{sub:vorh_sonnenwind}

Die Materieentstehungsrate im fraktalen Raum skaliert mit \( \dot{\rho}_{\text{cre}}(r) \propto r^{D-3} = r^{-0,296} \) (siehe Anhang~\ref{app:materieentstehung}). Diese Skalierung sollte sich in den Fluktuationen des Sonnenwinds widerspiegeln. Konkret:

\[
\delta v(r) \propto r^{-0,148}, \quad \delta \rho(r) \propto r^{-0,296}
\]

Diese Vorhersage ist mit bestehenden Satellitendaten (z. B. Parker Solar Probe, Solar Orbiter) testbar.

\subsection{Gravitationswellen-Nachglühen}
\label{sub:vorh_nachgluehen}

Bei Verschmelzungen von BEC-Objekten (im Gegensatz zu Schwarzen Löchern der ART) ist ein Nachglühen zu erwarten, da die Materie nicht hinter einem Ereignishorizont verschwindet (siehe Kapitel~\ref{ch:bec_objekte}). Dieses Nachglühen könnte im elektromagnetischen Bereich (Röntgen, Gamma) oder durch modulierte Gravitationswellen nachweisbar sein. Zukünftige Detektoren wie LISA oder das Einstein-Teleskop könnten diese Signatur von ART-Vorhersagen unterscheiden.

\subsection{Stromdichten in Plasmen}
\label{sub:vorh_stromdichten}

In Laborplasmen und astrophysikalischen Plasmen sollten Stromdichten mit \( j(r) \propto r^{D-3} \approx r^{-0,296} \) skalieren.

\section{Zusammenfassung der wichtigsten Vorhersagen}
\label{sec:vorh_tabelle_zusammenfassung}

Die folgende Tabelle fasst die wichtigsten testbaren Vorhersagen zusammen:

\begin{table}[htbp]
\centering
\small
\begin{tabularx}{\textwidth}{|l|X|X|}
\hline
\textbf{Vorhersage} & \textbf{Wert / Skalierung} & \textbf{Testmethode} \\
\hline
Galaxien-Dichteprofil & \(\rho(r) \propto r^{-0,296}\) & EUCLID, Rubin (LSST) \\
Rotationskurven & \(v(r) \propto r^{0,352}\) & Extrem breite Galaxien \\
CMB-Korrelationen & \(C(\theta) \propto \theta^{-0,296}\) & Planck, LiteBIRD \\
Neutrinomasse (Gen. 1) & \(m_{\nu,1} \approx 0,101 \, \text{eV}/c^2\) & KATRIN, Doppelbeta-Zerfall \\
Neutrinomasse (Gen. 2) & \(m_{\nu,2} \approx 9,3 \times 10^{-4} \, \text{eV}/c^2\) & Neutrino-Oszillationen \\
Neutrinomasse (Gen. 3) & \(m_{\nu,3} \ll 10^{-4} \, \text{eV}/c^2\) & Neutrino-Oszillationen \\
Oszillationslänge & \(L_{\text{osc}} \propto E^{0,775}\) & JUNO, DUNE, Hyper-K \\
Sonnenmasse & Kein Masseverlust durch Wind & Präzise Sonnenbeobachtung \\
Gravitationswellen-Dispersion & \(v_{\text{phase}}(\omega) = c(1 + \alpha \omega^{-0,099})\) & LISA, Einstein-Teleskop \\
Skalierung im Plasma & \(j(r) \propto r^{-0,296}\) & Laborplasmen, Sonnenwind \\
Lichtablenkung & \(\Delta\phi \propto \lambda^2\) & Multiband-Beobachtungen \\
Shapiro-Laufzeit & Frequenzabhängig & Pulsar-Timing \\
Maximale Kernmasse & \(M_{\text{BEC}} \approx 3 \times 10^6 M_\odot\) & Beobachtung von Galaxienzentren \\
EHT-Schatten & \(r_{\text{ph}} = 3GM/c^2\), andere Feinstruktur & EHT, nächste Generation \\
Hubble-Spannung & \(H_{\text{eff}}(d) \propto d^{0,704}\) & Supernovae, Quasare, GW \\
\hline
\end{tabularx}
\caption{Zusammenfassung der wichtigsten testbaren Vorhersagen der WDBT+.}
\label{tab:vorhersagen}
\end{table}

\section{Fazit}
\label{sec:vorh_fazit}

Die WDBT+ macht spezifische, testbare Vorhersagen, die sie von der Allgemeinen Relativitätstheorie und dem Standardmodell der Kosmologie unterscheiden. Besonders bemerkenswert ist die konsistente Erklärung zweier scheinbar unabhängiger Phänomene:

\begin{itemize}
    \item Die Neutrinomasse \( m_\nu \approx 0,1 \, \text{eV}/c^2 \)
    \item Die Hubble-Spannung \( H_{\text{eff}}(d) \propto d^{0,704} \)
\end{itemize}

Beide folgen aus derselben fraktalen Raumdimension \( D \approx 2,704 \) und derselben Korrelationslänge \( L_{Q,0} \approx 2{,}0 \times 10^{46} \, \text{m} \). Dies ist ein starkes Indiz für die Konsistenz der Theorie.

Sollten diese Vorhersagen bestätigt werden, wäre dies ein starkes Indiz für die Richtigkeit der Theorie. Sollten sie widerlegt werden, wäre die Theorie falsifiziert – ein Zeichen ihrer wissenschaftlichen Qualität.